\documentclass[11pt]{article}
\usepackage[a4paper,margin=1in]{geometry}
\usepackage[T1]{fontenc}
\usepackage[utf8]{inputenc}
\usepackage{lmodern}
\usepackage{hyperref}
\usepackage{xcolor}
\usepackage{listings}
\usepackage{longtable}

\lstdefinestyle{code}{
    basicstyle=\ttfamily\small,
    breaklines=true,
    frame=single,
    columns=fullflexible
}

\title{CircleNet Analytics Jobs Report\\Simple and Optimized Solutions (Task A--H)}
\author{DS503 Project 1}
\date{\today}

\begin{document}
    \maketitle

    \section{How We Approached the Work}
    For every task, we tried to do two versions:
    \begin{itemize}
        \item \textbf{Simple version}: direct and easy to understand.
        \item \textbf{Optimized version}: reduce extra shuffle/IO when possible.
    \end{itemize}

    We also added execution timing in each class using a small helper (\texttt{JobTimer}) so we can compare both versions using \texttt{total\_time\_ms}.

    \begin{lstlisting}[style=code,language=Java,caption={Timing helper used in all jobs}]
boolean ok = JobTimer.run(job, "TaskAOptimized");
JobTimer.total("TaskAOptimized", totalStart);
    \end{lstlisting}

    \section{Task A: Frequency of Favorite Hobby}
    \subsection*{Simple Solution}
    Read CircleNetPage, emit \texttt{(FavoriteHobby, 1)} in mapper, and sum in reducer.
    This is classic counting.

    \begin{lstlisting}[style=code,language=Java]
// Mapper output
context.write(new Text(hobby), new IntWritable(1));
    \end{lstlisting}

    \subsection*{Optimization Thinking and Implementation}
    This task is safe for combiner because sum is associative and commutative.
    So we used the same sum reducer as combiner to reduce network shuffle.

    \begin{lstlisting}[style=code,language=Java]
job.setCombinerClass(SumReducer.class);
    \end{lstlisting}

    \section{Task B: Top 10 Most Accessed Pages (Id, NickName, JobTitle)}
    \subsection*{Simple Solution}
    We used 3 jobs:
    \begin{itemize}
        \item Count accesses per page from ActivityLog.
        \item Join counts with CircleNetPage to get nickname/job title.
        \item Global top-10 selection.
    \end{itemize}

    \subsection*{Optimization Thinking and Implementation}
    We reduced one stage by loading CircleNetPage as a cache file in the top-10 job.
    So optimized flow is:
    \begin{itemize}
        \item Job 1: count page accesses with combiner.
        \item Job 2: single reducer keeps top-10 and map-side lookup from cache.
    \end{itemize}

    \begin{lstlisting}[style=code,language=Java]
job2.addCacheFile(new Path(args[1]).toUri()); // CircleNetPage.csv
    \end{lstlisting}

    \section{Task C: Users Whose Hobby Matches Chosen Hobby}
    \subsection*{Simple Solution}
    We pass hobby as argument (for example, \texttt{Chess}).
    Mapper checks \texttt{FavoriteHobby == target}, outputs \texttt{NickName, JobTitle}.
    Reducer only passes values.

    \subsection*{Optimization Thinking and Implementation}
    Reducer is not needed because there is no aggregation.
    So optimized version is map-only (\texttt{setNumReduceTasks(0)}), which saves IO.

    \begin{lstlisting}[style=code,language=Java]
job.setNumReduceTasks(0);
    \end{lstlisting}

    \section{Task D: Popularity Factor (Follower Count per Page Owner)}
    \subsection*{Simple Solution}
    We reduce-side joined:
    \begin{itemize}
        \item Page data: \texttt{(ownerId, P,nick)}
        \item Follows data: \texttt{(ownerId, F,1)} where ownerId is ID2
    \end{itemize}
    Reducer computes follower count and outputs zero if no follower rows were present.

    \subsection*{Optimization Thinking and Implementation}
    Optimized version first counts followers by ID2 (with combiner), then uses map-side join:
    load count output through Distributed Cache and scan CircleNetPage once.
    This avoids a full reduce-side join on all records.

    \section{Task E: Favorites Behavior (Total Actions + Distinct Pages Accessed)}
    \subsection*{Simple Solution}
    Mapper emits \texttt{(ByWho, WhatPage)}.
    Reducer counts:
    \begin{itemize}
        \item total actions = number of rows
        \item distinct pages = set size of \texttt{WhatPage}
    \end{itemize}

    \subsection*{Optimization Thinking and Implementation}
    Distinct counting in one reducer can be heavy.
    So optimized approach:
    \begin{itemize}
        \item Job 1 deduplicates \texttt{(ByWho, WhatPage)} pairs.
        \item Job 2 combines total-count stream and distinct-count stream.
    \end{itemize}
    This separates concerns and allows local aggregation with combiner in final stats stage.

    \section{Task F: Users More Popular Than Average}
    \subsection*{Simple Solution}
    We did it in 3 stages:
    \begin{itemize}
        \item Join pages with follower counts (including zero).
        \item Compute global average follower count.
        \item Filter users where follower count \(>\) average.
    \end{itemize}

    \subsection*{Optimization Thinking and Implementation}
    Optimized version:
    \begin{itemize}
        \item Count followers with combiner.
        \item Compute average once.
        \item Filter by scanning pages and doing cache lookup for follower count.
    \end{itemize}
    This avoids a larger reduce-side join output stage.

    \section{Task G: Outdated Pages (No Access in Last 90 Days)}
    \subsection*{Simple Solution}
    Based on dataset timestamp granularity (hour), 90 days = 2160 hours.
    Simple flow:
    \begin{itemize}
        \item Job 1 gets global max ActionTime.
        \item Job 2 gets each user's latest ActionTime.
        \item Job 3 joins with pages and keeps users with no activity or activity older than threshold.
    \end{itemize}

    \subsection*{Optimization Thinking and Implementation}
    Optimized version removes one MR stage:
    \begin{itemize}
        \item Job 1 computes user latest ActionTime (combiner = max).
        \item Driver computes global max by reading job output.
        \item Job 2 map-side joins with pages using cache and filters outdated directly.
    \end{itemize}

    \section{Task H: Same Region One-Way Follows (Not Followed Back)}
    \subsection*{Simple Solution}
    Steps:
    \begin{itemize}
        \item Load page regions using cache.
        \item Keep only follow edges where both users are in same region.
        \item For each unordered pair, check if relation is one-way only.
        \item Join result with pages to output \texttt{ID, nickname}.
    \end{itemize}

    \subsection*{Optimization Thinking and Implementation}
    Optimized approach moves nickname lookup into reducer (cache-based), then runs a light dedup job.
    So we avoid one larger reduce-side join stage.

    \section{How We Compare Simple vs Optimized}
    For each run, logs print per-job and total time:
    \begin{itemize}
        \item \texttt{... time\_ms=...}
        \item \texttt{... total\_time\_ms=...}
    \end{itemize}
    For comparison table, we use \texttt{total\_time\_ms}.

    \begin{longtable}{|l|r|r|}
        \hline
        \textbf{Task} & \textbf{Simple (ms)} & \textbf{Optimized (ms)}\\
        \hline
        A &  & \\
        \hline
        B &  & \\
        \hline
        C &  & \\
        \hline
        D &  & \\
        \hline
        E &  & \\
        \hline
        F &  & \\
        \hline
        G &  & \\
        \hline
        H &  & \\
        \hline
    \end{longtable}

    \section{Short Reflection}
    As a student team, we first wrote clear solutions that we can explain easily.
    Then we optimized carefully using simple Hadoop ideas:
    combiner where safe, map-only where possible, and map-side cache joins for smaller reference data.
    This gave us better runtime while still keeping code understandable.

\end{document}
